% !TEX root =  main.tex
\chapter{The Detection Challenge}
\label{chap:2}

\section{From Pollen Grains}
\label{sec:2.1}
\hspace{\parindent} The Brownian Motion story begins with a mystery no one could explain. In 1827, the botanist Robert Brown watched pollen grains suspended in water twitch under his microscope with no apparent cause, and reported the observation without being able to account for it \cite{Brown1828}. For nearly seventy-eight years the jiggling carried his name but no accepted mechanism. Albert Einstein supplied one explanation in 1905, arguing ceaselessly battering the suspended grain from random directions \cite{Einstein1905, Einstein1956}. Marian Smoluchowski arrived at an equivalent picture independently around the same time \cite{Smoluchowski1906}. What made Einstein's account more than a plausible story was that it made a number falsifiable: it predicted how far, on average, a particle of known size should wander in a fluid of known viscosity and temperature. Jean Perrin tested that prediction directly under the microscope a few years later, and the resulting agreement was treated as some of the strongest available evidence that atoms and molecules were physically real rather than a convenient bookkeeping device \cite{Perrin1916}.

\hspace{\parindent} Paul Langevin gave the same physics a different, and for simulation purposes more usable form. Rather than track a probability distribution, he wrote down a force equation for a single particle: a drag term that always opposes motion, plus a rapidly fluctuating random force representing the unresolved molecular collisions \cite{Coffey2004}. This is the Langevin equation already introduced in Section~\ref{sec:1.2}, and it is the equation LAMMPS integrates, through its \texttt{fix langevin} command, for every probe and foreign particle in this thesis' simulations. Ornstein and Uhlenbeck later worked out the statistical consequences of that equation in full: over sampling intervals much longer than the particle's momentum relaxation time, displacements settle into a zero-mean Gaussian distribution \cite{Uhlenbeck1930}. That single result is what makes two independent measurement strategies possible later in this thesis where one can either measure how far particles spread out on average through the mean squared displacement route or measure the shape of the displacement distribution itself through the use of the variance route, and both should recover the same diffusion coefficient if the motion is truly Brownian.

\section{Two ways to Measure the Same Number}
\label{sec:2.2}
\hspace{\parindent} Knowing that Brownian motion should be Gaussian and diffusive is not the same as being able to extract a clean diffusion coefficient from a finite, noisy simulated or experimental trajectory. This is where the Brownian group's own prior work becomes directly irrelevant. Cusi, Tapang, and Saloma first examined how probe-probe collisions corrupt diffusion coefficient measurements in dilute suspensions \cite{Cusi2018}, then followed up with a systematic comparison of the two extraction methods: reading the diffusion coefficient off the slope of the mean squared displacement (MSD) curve, $D_\text{MSD}$, versus reading it off the variance of the displacement distribution, $D_\text{var}$ \cite{Cusi2020}. Cusi's undergraduate thesis extended this further into an experimental single-particle tracking study, and identified spatial resolution, which is how finely a camera can resolve a particle's position, as the single most influential factor governing whether a measured diffusion coefficient can be trusted \cite{Cusi2020Thesis}. Around the same period, Michalet and Berglund arrived at a related conclusion from a more theoretical direction, showing that the accuracy of diffusion coefficient estimates from tracked trajectories depends sensitively on how the fit is constructed, not just on how good the raw data are \cite{Michalet2012}.

\hspace{\parindent} The immediate predecessor to this thesis, Balejado, Escaran, Cusi, and Saloma's 2025 conference paper, asked a narrower and more pointed questions: given that $D_\text{MSD}$ and $D_\text{var}$ should agree in principle, at what sampling time interval $\tau$ do they actually agree in LAMMPS simulation \cite{balejado2025} The answer hinges on the Langevin damping parameter $\tau_\text{damp}$, the characteristic time over which a particle forgets its own momentum through viscous drag. Sampling faster than $\tau_\text{damp}$ catches the particle still moving ballistically along its last random kick, which biases both methods low; sampling much slower than $\tau_\text{damp}$ recovers the correct diffusive value but shortens the usable trajectory and inflates statistical noise. Related diagnostic work along the same lines, focused specifically on the reliability of computed diffusion coefficients and their uncertainties, was carried out by McCluskey, Coles, and Morgan \cite{McCluskey2023}. The practical result for this thesis is a design rule rather than a debate: choose $\tau \gg \tau_\text{damp}$ before trusting a single measurement, and confirm the simulation is behaving as pure Brownian motion before asking any question about detection. Chapter~\ref{chap:3} follows exactly that rule.


\section{The Detection Problem: How Long Until Something is Found?}
\label{sec:2.3}
Diffusion theory describes how a particle spreads out; it does not, by itself, say when that particle will first reach a specific target. That question belongs to a distinct body of work known as first-passage or search theory, laid out comprehensively by Redner \cite{Redner2001} and extended to biological chemical search problems by B\'enichou, Loverdo, Moreau, and Voituriez \cite{Loverdo2008, Benichou2011}. Their central result, already stated as Equation~\ref{eqn:1.1}, is deceptively simple where if a lone searcher takes an average time $\langle T \rangle$ to find a target, then $N_p$ independent, non-interacting searchers should find it in $\langle T \rangle / N_p$ on average. Condamin and colleagues showed that this scaling holds even in geometrically complicated, scale-invariant search spaces \cite{Condamin2007}, while B\'enichou and their collaborators demonstrated that the geometry of the search domain itself can be tuned to control how quickly a target is found \cite{Benichou2010}. Two features of this literature matter directly for this thesis. First, the $1/N_p$ law is explicitly an \emph{ideal-limit} statement, valid only while searchers genuinely do not interfere with one another; nothing in the theory says what happens once they do. Second, the law rests on ergodicity, the idea that averaging over many independent searchers observed briefly gives the same answer as averaging over one searcher observed for a proportionally longer time. The ergodic product, $E = N_p \times \langle T \rangle$, used throughout this thesis, is simply a way of putting that assumption to a direct test and if it holds, $E$ should sit flat against $N_p$, and any systematic drift is evidence that the searchers have stopped behaving independently.

\section{Crowding and Caging}
\label{sec:2.4}
\hspace{\parindent} The search theory literature is largely built for dilute systems, where adding a searcher costs nothing but its own trajectory. Real search environments rarely stay dilute once enough searchers are added, and a separate body of work in soft condensed matter physics describes what happens when they do not. Weeks, Crocker, Levitt, Schofield, and Weitz used three-dimensional microscopy of colloidal suspensions to directly visualize the onset of caging where as particle packing increases toward the colloidal glass transition, each particle becomes increasingly hemmed in by its neighbors, and its mobility drops sharply even though nothing has frozen in the conventional sense \cite{Weeks2000}. Torquato's treatment of random heterogeneous materials provides the geometric backbone for quantifying how densely particles can be packed before free space effectively disappears, including the hexagonal close packing bound $\phi_\text{HCP} = \pi/(2\sqrt{3}) \approx 0.9069$ used as a reference ceiling in Section~\ref{sec:1.2} \cite{Torquato2002}. The physical intuition these two works supply is what motivates the crowding-limited regime described in this thesis where past some area fraction $\phi$, newly added probes do not meaningfully explore new space, they mostly collide with, and get boxed in by, the probes and foreign particles that are already there. This is the same physical mechanism, expressed through the excluded-volume repulsion of the Weeks-Chandler-Andersen (WCA potential \cite{Weeks1971}, that ultimately produces the upturn in $\langle T \rangle$ at high $N_p$ documented in Balejado, Tapang, and Saloma's 2026 conference paper \cite{Balejado2026} and revisited at greater depth in this thesis. 

\section{LAMMPS}
\label{sec:2.5}
\hspace{\parindent} None of the above would be testable at scale without the reliable simulation engine. This thesis uses the Large-scale Atomic/Molecular Massively Parallel Simulator (LAMMPS), a particle-based molecular dynamics code originally developed for materials science but widely repurposed for soft-matter and colloidal simulation because of its flexible pairwise potentials and built-in Langevin thermostat features \cite{Thompson2022}. Visual inspection and rendering of the resulting particle trajectories, which is useful during the development and debugging phase, relied on OVITO, the Open Visualization Tool \cite{Stukowski2010}.

\hspace{\parindent} One further statistical tool deserved mention because it plays a direct role in Chapter~\ref{chap:4}'s presentation of results which is the violin plot. A conventional error bar communicates only a mean and a spread, and can look identical whether the underlying detection times are tightly clustered, evenly spread, or badly skewed with a long tail of rare, very slow detections. Hintze and Nelson introduced the violin plot precisely to solve this problem, combining a traditional box-plot summary with a mirrored kernel density trace so that the full shape of a distribution, not just its first two moments, is visible at a glance \cite{Hintze1998}. Given that this thesis specifically expects detection time distributions to change shape, from tight and symmetric in the diffusion-limited regime to broad and right-skewed once caging sets in, the violin plot is not a cosmetic choice but a diagnostic one, and it is treated as such in the results and discussion chapter.

\section{The Study Gap this Thesis Addresses}
\label{sec:2.6}
\hspace{\parindent} Taken together, the literature above supplied every ingredient except the one this thesis contributes. Search theory gives the ideal $1/N_p$ scaling law and the ergodic product as its diagnostic. Colloidal physics gives the caging mechanism that should eventually break that law. Our group's own prior work gives a validated LAMMPS pipeline and a settled, damping-parameter-aware protocol for extracting trustworthy diffusion coefficients. What had not yet been done, at the start of this project, was to put a size-varying, two-species, WCA-repulsive Brownian mixture through the full range of probe densities, from a single lonely searcher up to a near-close-packed swarm, and ask exactly where along that range the diffusion-limited and crowding-limited regimes meet, and whether the ergodic product can flag that transition earlier than tthe raw detection time can. That question, first posed in preliminary form in \cite{Balejado2026}, is the one that this thesis answers in full, across multiple systet sizes $N_\text{total}$, in Chapters~\ref{chap:3} and~\ref{chap:4}.